\uuid{fX3d}
\exo7id{3517}
\titre{exo7 3517}
\auteur{quercia}
\organisation{exo7}
\datecreate{2010-03-10}
\isIndication{false}
\isCorrection{true}
\chapitre{Réduction d'endomorphisme, polynôme annulateur}
\sousChapitre{Diagonalisation}
\cpge{mpsi-mp2i.changements-bases.c, mp-mpi.reduction-endomorphismes.d, pc.reduction.c, psi.reduction-endomorphismes.c, pt.reduction-endomorphismes.c, mp-mpi.reduction-endomorphismes.i, pc.reduction.d, psi.reduction-endomorphismes.d, bcpst2.valeurs-propres-vecteurs-propres.b}
\module{Algèbre}
\niveau{L2}
\difficulte{}

\contenu{
\question{
Soit $A\in\mathcal{M}_n(\C)$. L'endomorphisme $f$ de $\mathcal{M}_n(\C)$ défini par
$f(M) = \mathrm{tr}(A)M+\mathrm{tr}(M)A$ est-il diagonalisable~?
}
\reponse{
Oui ssi $\mathrm{tr}(A)\ne 0$ ou $A=0$.
}
}
